\scanpage{55}
4.2. Définition du site et du topos cristallins \inserted{de $X$ sur $S$}

Définition du site et du topos stratifiant de $X$ sur $S$

C'est pareil si $X$ form\textsuperscript{t}. lisse sur $S$.

Ce sont des topos annelés

\deleted{Description d'un} \inserted{Attention, le site n'\uline{a pas en général d'objet final} !}

Ce sont des \deleted{\unsure{sites|Sit} \unsure{fais}} et topos \uline{annelés}.

Description des \deleted{\unsure{un}} faisceaux d'ensembles en termes de faisceaux habituels. Idem pour faisceaux d'anneaux, faisceaux de Modules sur faisceaux d'anneaux etc.

Faisceaux spéciaux : les hom de changement de base sont des isom.

\diagram{\begin{tikzcd}[ampersand replacement=\&,row sep=small]
\text{Faisceaux de Modules spéciaux sur site stratifiant \inserted{de $X/S$}}
\arrow[d,phantom,"\rotatebox{90}{$\simeq$}" description] \\
\text{Faisceaux de Modules à stratification sur $X/S$.}
\end{tikzcd}}

4.3. Énonçons de suite les résultats et conjectures suivants, qui seront \deleted{\unsure{étud|étudiés}} détaillés plus tard :

\uline{Th} Si $S$ de car. $0$ et $X$ lisse sur $S$, alors
\[
H^*_{\substack{\mathrm{CRIS}\deleted{\unsure{/S|X/S}}\\\text{ou }\mathrm{STRAT}}}(X/S,\underline O_X)\simeq H^*_{\mathrm{DR}}(X/S)
\]
(isom can). \deleted{et en fait $R f_*(\underline O_X)$ est \uline{spécial}, et
\[
(R f_{\unsure{*\mathrm{cris}|*}}(\underline O_X))^{\unsure{\mathrm{ord}}}\simeq\mathbb R f_*(\Omega^{\cdot}_{X/S})\quad ....
\]}
% archive sheet number: 54

